
{\actuality} 

Метод молекулярной динамики прочно утвердился в качестве мощного вычислительного инструмента для исследования широкого спектра задач, требующих атомистического разрешения. Уникальная сила МД заключается в его способности комплементарно заменять или существенно дополнять сложные и дорогостоящие эксперименты в ситуациях, когда прямой контроль над процессами невозможен или сопряжен с большими трудностями. Это особенно актуально для изучения экстремальных условий, таких как сверхвысокие давления или температуры, воссоздание которых в лаборатории требует уникального оборудования, или исследования быстропротекающих процессов, например, начальных стадий нуклеации или механизмов химических реакций, где метод молекулярной динамики предоставляет временное разрешение, недостижимое для большинства экспериментальных методик .

Кроме того, для многих сложных систем — будь то молекулярные жидкости, расплавы или наноразмерные объекты — метод молекулярной динамики предоставляет детализированное, на молекулярном уровне, понимание происходящих процессов. В отличие от методов, дающих усредненную картину, метод молекулярной динамики позволяет непосредственно наблюдать и анализировать атомные перемещения, формирование и разрушение связей, что является ключом к установлению связей между структурой, динамикой и макроскопическими свойствами вещества .

Как и любой метод, метод молекулярной динамики имеет ряд принципиальных ограничений, главным из которых является размер исследуемой системы. Существует неизбежный компромисс между точностью модели и ее вычислительной стоимостью: более точные модели, такие как ab initio метод молекулярной динамики, требуют колоссальных ресурсов, что ограничивает размер системы и время моделирования. Напротив, грубые модели, использующие классические силовые поля, позволяют имитировать поведение миллионов атомов на наносекундных масштабах, но могут проигрывать в точности и предсказательной силе, особенно для свойств, критически зависящих от электронной структуры.

Данная работа направлена на поиск практического компромисса между стремлением к точности и увеличением размера системы. Проводится систематическое сравнение результатов, полученных с помощью потенциалов межатомного взаимодействия различной сложности для ряда физически значимых систем. Центральная гипотеза исследования состоит в том, что, даже если абсолютные значения величин, предсказываемые разными моделями, расходятся, качественный вид функциональных зависимостей может воспроизводиться удовлетворительно. Это открывает возможность калибровки более простых и вычислительно эффективных моделей по данным сложных расчетов или ограниченному набору экспериментальных точек с последующей экстраполяцией результатов в широком диапазоне условий.

Для проверки этой гипотезы было рассмотрено несколько систем и процессов. Первый из них \--- самодиффузия в углеводородах: на примере 2,2,4-триметилпентана ($C_8H_{18}$) и 1,1-дифенилэтана ($C_{14}H_{14}$) исследовано влияние выбранного потенциала межатомного взаимодействия на расчет плотности и коэффициента самодиффузии.

Поверхностные явления в двухфазной системе рассмотрены на примере смеси метана и этана в двухфазной области. Рассчитано поверхностное натяжение и получена величина парахора. Показано, что модель объединенного атома позволяет воспроизвести это значение, что дает метод прогноза поверхностного натяжения с достаточной точностью без проведения экспериментов в целевых условиях.

Адсорбция на межфазной границе была рассмотрена для той же бинарной смеси. Выведена формула для изотермы адсорбции, не использующая приближений идеального газа и идеального раствора, что повышает ее применимость для реальных систем.

На примере наночастиц железа методом метод молекулярной динамики исследованы элементарные акты соударения единичных атомов с кластером. Эти процессы — термическая аккомодация и прилипание — являются ключевыми для понимания кинетики роста наночастиц в процессе нуклеации и конденсации из пара . Именно они в значительной степени определяют конечное распределение кластеров по размерам, которое, в свою очередь, диктует их каталитические, оптические и магнитные свойства.

% {\progress}
% Этот раздел должен быть отдельным структурным элементом по
% ГОСТ, но он, как правило, включается в описание актуальности
% темы. Нужен он отдельным структурынм элемементом или нет ---
% смотрите другие диссертации вашего совета, скорее всего не нужен.

{\aim} работы \--- расширение возможностей метода молекулярной динамики в прогнозировании термодинамических, транспортных и поверхностных свойств сложных систем путём разработки и верификации методик на основе потенциалов межатомного взаимодействия различной сложности.

Для достижения поставленной цели необходимо было решить следующие {\tasks}.
\begin{enumerate}[beginpenalty=10000]
  \item Разработана и апробирована методика сравнительного анализа потенциалов межатомного взаимодействия на примере расчёта плотности и коэффициента самодиффузии 2,2,4-триметилпентана ($C_8H_{18}$) и 1,1-дифенилэтана ($C_{14}H_{14}$) в широком диапазоне давнений (0.1–400 МПа). 
  \item Разработан подход к прогнозированию поверхностного натяжения на основе расчёта парахора методом молекулярной динамики. Методика применена к системе «пар–жидкость» чистого этана и бинарной смеси метан–этан при варьируемых температуре и давлении.
  \item Предложена и верифицирована модификация уравнения изотермы адсорбции метана на поверхности жидкой фазы этана, исключающая приближения идеального газа и идеального раствора. Адекватность уравнения подтверждена прямым сравнением с результатами молекулярно-динамического моделирования.
  \item Создана методика расчёта коэффициентов термической аккомодации и прилипания для наноразмерных кластеров железа, позволяющая исследовать зависимость этих параметров от температуры (200–2500 К) и размера частицы.
\end{enumerate}

{\novelty}
\begin{enumerate}[beginpenalty=10000] % https://tex.stackexchange.com/a/476052/104425
  \item Впервые предложена аналитическая формула для описания адсорбции на межфазной границе в двухкомпонентной системе, не использующая приближений идеального газа и идеального раствора. Данный подход позволяет с повышенной точностью моделировать равновесные свойства границы раздела фаз в реальных системах.
  \item Разработан новый метод прогнозирования поверхностного натяжения, основанный на комбинации одного экспериментального измерения (плотности фаз и поверхностного натяжения в реперной точке) с расчётом парахора методом молекулярной динамики в рамках модели объединённого атома. Метод позволяет предсказывать свойства в широком диапазоне параметров состояния без проведения дополнительных экспериментов.
  \item Установлены фундаментальные закономерности в поведении коэффициента термической аккомодации для наноразмерных кластеров железа. Впервые показано, что зависимость коэффициента от температуры кластера имеет аррениусовский характер, а от размера кластера — масштабируется как $N^{-1/3}$, что свидетельствует о ключевой роли поверхностных эффектов.
  \item Проведено систематическое исследование влияния выбора потенциала межатомного взаимодействия на предсказание плотности и коэффициента самодиффузии для 2,2,4-триметилпентана ($C_8H_{18}$) и 1,1-дифенилэтана ($C_{14}H_{14}$) в диапазоне давлений 0.1–400 МПа. Количественно оценены границы применимости различных классов потенциалов, что имеет ключевое значение для планирования вычислительных экспериментов.
\end{enumerate}

{\influence} Разработана фундаментальная модель адсорбции. Впервые предложена аналитическая формула для изотермы адсорбции на межфазной границе, не использующая приближений идеального газа и идеального раствора. Данный подход позволяет количественно описывать адсорбционные явления в широком диапазоне концентраций и условий, выходя за рамки классической теории, ограниченной областью малых концентраций.

Установлены ключевые закономерности процессов на поверхности нанокластеров. Подтверждена гипотеза о степенной зависимости коэффициента термической аккомодации от размера кластера ($\sim N^{-1/3}$) и её аррениусовском характере от температуры. Полученные результаты вносят существенный вклад в понимание атомистических механизмов тепло- и массопереноса на начальных стадиях нуклеации, предоставляя количественную основу для разработки более точных кинетических моделей роста наночастиц.

Разработан эффективный метод прогнозирования поверхностных свойств. Показано, что расчёт парахора методом молекулярной динамики в рамках вычислительно эффективной модели объединённого атома демонстрирует удовлетворительное согласие с экспериментальными данными. Это позволяет, используя единственное реперное измерение плотностей фаз и поверхностного натяжения, прогнозировать всю температурную зависимость поверхностного натяжения, что существенно сокращает потребность в дорогостоящих и трудоёмких экспериментах.

{\methods} Молекулярно-динамическое моделирование выполнялось с использованием программного пакета LAMMPS (Large-scale Atomic/Molecular Massively Parallel Simulator). Расчеты проводились на суперкомпьютерных системах «DESMOS» и «Fisher».

Для анализа транспортных свойств коэффициент самодиффузии рассчитывался в рамках равновесного подхода по методу Эйнштейна-Смолуховского через анализ зависимости среднеквадратичного смещения частиц от времени.

Исследование поверхностных явлений проводилось методом двухфазного моделирования с явным учётом сосуществования жидкой и паровой фаз в равновесных условиях. Величина адсорбции на межфазной границе определялась в рамках гиббсовского формализма через анализ профилей концентраций компонентов.

Коэффициенты термической аккомодации и прилипания рассчитывались по методике, описанной в работе~\cite{insepov_kinetics_1991}, путём статистического анализа изменения энергии и импульса атомов при столкновении с поверхностью кластера.

Визуализация атомных конфигураций, траекторий и анализ структуры проводились с помощью программного пакета OVITO (Open Visualization Tool)~\cite{ovito}.

{\defpositions}
\begin{enumerate}[beginpenalty=10000] % https://tex.stackexchange.com/a/476052/104425

  \item Несмотря на несовпадение зависимостей плотности и коэфффициента самодиффузии от давления с экспериментальными данными, коррекция плотностей позволит получить значения коэффициента диффузии, более близкие к референтным.
  \item Воспроизведено значение парахора чистого этана для давлений 4-40 атм, что позволит использовать метод молекулярной динамики для предсказания экспериментальных значений поверхностного натяжения.
  \item Формула для оценки адсорбции на плоской поверхности Гиббса для двухфазной двухкомпонентной системы, использующая только экспериментально измеряемые величины без приближений идеального газа и идеального раствора, проверена для двухфазной смеси, состоящей из метана и этана.
  \item Методика расчета коэффициента термической аккомодации атомов окружающего газа на малых кластерах в рамках метода молекулярной динамики, получены его зависимости от температуры кластера и его размера для кластеров железа и налетающих на них атомов аргона.
  \item Количество уносимой от кластера налетающим атомом энергии зависит от температуры кластера слабее, чем по линейному закону, что объясняется уменьшением времени взаимодействия атома с кластером при повышении температуры.
  
\end{enumerate}

{\reliability} Результаты молекулярно-динамических расчетов находятся в согласии с известными экспериментальными данными.

{\probation}
Основные результаты работы были представлены на конференциях МФТИ(г. Долгопрудный) в 2014-2020, 2023 и 2024 годах, на международных конференциях International Conference on Equations of State for Matter(пос. Эльбрус) в 2016, 2020 и 2022 годах, International Conference on Equations of State for Matter(пос. Эльбрус) в 2017 и 2019 годах, Российские симпозиумы «Фундаментальные основы атомистического многомасштабного моделирования»(г. Новый Афон) в 2015-2019 годах, The XXth Research Workshop Nucleation Theory and Applications(г. Дубна) в 2016 году.

{\contribution} Все представленные в диссертации результаты получены лично автором. Постановка задач, вошедших в диссертационную работу, выполнена под руководством В.В. Писарева и Г.Э. Нормана. Выводы и положения, выносимые на защиту, сформулированы лично автором.

\ifnumequal{\value{bibliosel}}{0}
{%%% Встроенная реализация с загрузкой файла через движок bibtex8. (При желании, внутри можно использовать обычные ссылки, наподобие `\cite{vakbib1,vakbib2}`).
    {\publications} Основные результаты по теме диссертации изложены
    в~XX~печатных изданиях,
    X из которых изданы в журналах, рекомендованных ВАК,
    X "--- в тезисах докладов.
    
}%
{%%% Реализация пакетом biblatex через движок biber
    \begin{refsection}[bl-author, bl-registered]
        % Это refsection=1.
        % Процитированные здесь работы:
        %  * подсчитываются, для автоматического составления фразы "Основные результаты ..."
        %  * попадают в авторскую библиографию, при usefootcite==0 и стиле `\insertbiblioauthor` или `\insertbiblioauthorgrouped`
        %  * нумеруются там в зависимости от порядка команд `\printbibliography` в этом разделе.
        %  * при использовании `\insertbiblioauthorgrouped`, порядок команд `\printbibliography` в нём должен быть тем же (см. biblio/biblatex.tex)
        %
        % Невидимый библиографический список для подсчёта количества публикаций:
        \printbibliography[heading=nobibheading, section=1, env=countauthorvak,          keyword=biblioauthorvak]%
        \printbibliography[heading=nobibheading, section=1, env=countauthorwos,          keyword=biblioauthorwos]%
        \printbibliography[heading=nobibheading, section=1, env=countauthorscopus,       keyword=biblioauthorscopus]%
        \printbibliography[heading=nobibheading, section=1, env=countauthorconf,         keyword=biblioauthorconf]%
        \printbibliography[heading=nobibheading, section=1, env=countauthorother,        keyword=biblioauthorother]%
        \printbibliography[heading=nobibheading, section=1, env=countregistered,         keyword=biblioregistered]%
        \printbibliography[heading=nobibheading, section=1, env=countauthorpatent,       keyword=biblioauthorpatent]%
        \printbibliography[heading=nobibheading, section=1, env=countauthorprogram,      keyword=biblioauthorprogram]%
        \printbibliography[heading=nobibheading, section=1, env=countauthor,             keyword=biblioauthor]%
        \printbibliography[heading=nobibheading, section=1, env=countauthorvakscopuswos, filter=vakscopuswos]%
        \printbibliography[heading=nobibheading, section=1, env=countauthorscopuswos,    filter=scopuswos]%
        %
        \nocite{*}%
        %
        {\publications} Основные результаты по теме диссертации изложены в~\arabic{citeauthor}~печатных изданиях,
        \arabic{citeauthorvak} из которых изданы в журналах, рекомендованных ВАК\sloppy%
        \ifnum \value{citeauthorscopuswos}>0%
            , \arabic{citeauthorscopuswos} "--- в~периодических научных журналах, индексируемых Web of~Science и Scopus\sloppy%
        \fi%
        \ifnum \value{citeauthorconf}>0%
            , \arabic{citeauthorconf} "--- в~тезисах докладов.
        \else%
            .
        \fi%
        \ifnum \value{citeregistered}=1%
            \ifnum \value{citeauthorpatent}=1%
                Зарегистрирован \arabic{citeauthorpatent} патент.
            \fi%
            \ifnum \value{citeauthorprogram}=1%
                Зарегистрирована \arabic{citeauthorprogram} программа для ЭВМ.
            \fi%
        \fi%
        \ifnum \value{citeregistered}>1%
            Зарегистрированы\ %
            \ifnum \value{citeauthorpatent}>0%
            \formbytotal{citeauthorpatent}{патент}{}{а}{}\sloppy%
            \ifnum \value{citeauthorprogram}=0 . \else \ и~\fi%
            \fi%
            \ifnum \value{citeauthorprogram}>0%
            \formbytotal{citeauthorprogram}{программ}{а}{ы}{} для ЭВМ.
            \fi%
        \fi%
        % К публикациям, в которых излагаются основные научные результаты диссертации на соискание учёной
        % степени, в рецензируемых изданиях приравниваются патенты на изобретения, патенты (свидетельства) на
        % полезную модель, патенты на промышленный образец, патенты на селекционные достижения, свидетельства
        % на программу для электронных вычислительных машин, базу данных, топологию интегральных микросхем,
        % зарегистрированные в установленном порядке.(в ред. Постановления Правительства РФ от 21.04.2016 N 335)
    \end{refsection}%
    \begin{refsection}[bl-author, bl-registered]
        % Это refsection=2.
        % Процитированные здесь работы:
        %  * попадают в авторскую библиографию, при usefootcite==0 и стиле `\insertbiblioauthorimportant`.
        %  * ни на что не влияют в противном случае
        \nocite{vakbib2}%vak
        \nocite{patbib1}%patent
        \nocite{progbib1}%program
        \nocite{bib1}%other
        \nocite{confbib1}%conf
    \end{refsection}%
        %
        % Всё, что вне этих двух refsection, это refsection=0,
        %  * для диссертации - это нормальные ссылки, попадающие в обычную библиографию
        %  * для автореферата:
        %     * при usefootcite==0, ссылка корректно сработает только для источника из `external.bib`. Для своих работ --- напечатает "[0]" (и даже Warning не вылезет).
        %     * при usefootcite==1, ссылка сработает нормально. В авторской библиографии будут только процитированные в refsection=0 работы.
}

Ленев Д.Ю., Захаров С.А., Писарев В.В., Поверхностное натяжение и адсорбция на границе пар–жидкость в системе метан–этан , ЖФХ, 2024, 98, 12, 127-133. 

     N. Kondratyuk, D. Lenev, V. Pisarev, Transport coefficients of model lubricants up to 400 MPa from molecular dynamics, J CHEM PHYS, 2020, 152, 19, 191104

Д.Ю. Ленёв, Г.Э. Норман, Молекулярное моделирование термической аккомодации атомов аргона на кластерах атомов железа , HIGH TEMP+, 2019, 57, 4, 534–542

D. Yu. Lenev, G. E. Norman, Sticking coefficient for Fe atoms interacting with iron cluster, Journal of Physics: Conference Series, 2018, 946, 1, 012111

D.Yu. Lenev, G.E. Norman, Thermal accommodation at cold Ar atoms collisions with small Fe clusters at different temperatures , Journal of Physics: Conference Series, 2016, 774, 1, 6